\section{Poves kvadratne opne}

Poissonovo enačbo
\begin{equation}
    \nabla^2 u(x,y) = g(x,y)
\end{equation}
rešujemo direktno s Fourierovo transformacijo. Zaradi Dirichletovih robnih pogojev ($u_{RP}=0$) je ceneje uporabiti hitro (diskretno) sinusno transformacijo (DFT) v dveh dimenzijah, kjer z malimi črkami označujemo veličine v časovnem in z velikim v frekvenčnem prostoru
\begin{equation}
    G^{m n} = \sum_{j=1}^{J-1} \sum_{\ell=1}^{L-1} g_{j \ell} \, \sin  \frac{\pi m j}{J} \sin \frac{\pi n \ell}{L}.
\end{equation}
Alternativa je dvojna DFT v eni dimenziji, posebej po vrsticah in stolpcih.
Tako se potem glasi rešitev v frekvenčnem prostoru
\begin{equation}
    U^{m n} = \frac{h^2}{2}\,\frac{G^{m n}}{2\,\cos{\frac{\pi m}{J}} + \cos{\frac{\pi n}{L}} -2},
\end{equation}
kjer je $h=h_x=h_y$ diskretizacijski korak. Nazaj v fizikalni prostor se vrnemo z inverzno DFT v 2D.

Problem lahko rešimo tudi s kombinirano metodo, kjer izvedemo 1D DST, nam v preostali smeri ostane tridiagonalni sistem enačb za delne transformiranke $U_j^m$
\begin{equation}
    U_{\ell-1}^m -U_{\ell}^m\,(4-2\,\cos{\frac{\pi m}{J}}) + U_{\ell+1}^m = h^2\,G_{\ell}^m,
\end{equation}

Za drugi način po Metodi končnih razlik (MKR)  diskretiziramo računsko domeni in rešujemo dobljen sistem linearnih enačb eksaktno ali iterativno (SOR, Gaus-Seidel, Jacobi...), kot smo to storili v nalogi 6.

Uporabil sem algoritme za diskretno sinusno transformacijo (DST) iz Python knjižnice \texttt{scipy.fft}.

\subsection{Homogena opna}

V tem primeru je desna stran enačbe (porazdelitev mase) $g(x,y) = g_{j\ell} = 1$ po celotni opni. Za ta primer lahko uporabimo tudi naivno seštevanje Fourierovih vrst po enačbi, ki je definirana v navodilih, vendar, kot se izkaže, ta način ni preveč hiter; uporaba FFT algoritma za izračun diskretne sinusne transformacije (DST) precej pospeši računanje.

Na spodnjih slikah vidimo, da se rešitve na homogeni opni razlikujejo najbolj glede velikosti povesa na sredini opne. Opazim, da uporaba 1D oz. 2D DFT ne vpliva na rešitev. Kombinirana metoda presenetljivo vrne največji poves, vendar je mogoče le, da konvergenca poteka z druge strani! Rešitev naivnega seštevanja Fourierovih vrst, kar imenujem `zvezna' Fourierova metoda, se obnese podobno kot MKR. Ker ni opaziti večjih razlik med rešitvami iterativnih metod in eksaktne rešitve sistema enačb (to ni eksaktna rešitev problema!), lahko menim, da je dosežena dovoljšna natančnost. Pri tem opazim, da je SOR občutno najhitrejša.

\fig{H}{1}{./10.nal_slike/metode_homogena_r1}{Poves homogene opne, izračunan z različnimi metodami}{}
\fig{H}{1}{./10.nal_slike/prerez_homogena_r1}{Prerez po sredini preko koordinatnih osi}{}
\fig{H}{0.85}{./10.nal_slike/Umax_homogena_r1}{Največji poves homogene opne po različnih metodah}{}

Rešitev na homogeni opni poizkušamo aproksimirati s polinomom in znano vrednostjo največjega povesa $u_{\max}$ na srednini opne 
\begin{equation}
    u_p(u_{\max},x,y) = u_{\max}\;16\;(x-x^2)\,(y-y^2).
\end{equation}
Vendar se izkaže, da je tak popis podoben rešitvi, ni pa popoln!
\fig{H}{0.7}{./10.nal_slike/konvergenca/residuali_homog}{Odstopanje od polinomske rešitve}{} 

\subsection{Pravokotna obtežitev}

Preizkusim še pravokotno obtežitev na različnih mestih na opni, ki jo definiram kot 
\begin{equation}
    g(x,y) = \begin{cases}
        ratio, & \text{če } x\in \mathcal{A},\\
        1, & \text{sicer}.
    \end{cases}
\end{equation}
$\mathcal{A}$ opisuje črn pravokotnik na sliki, kjer je masa opne drugačna. Za primer na robu to opisujeta neenačbe $0.5<x<1$ in $0.5<y<1$, v primeru središčne lege pa $0.4<x<0.6$ in $0.4<y<0.6$. Razmerje med obtežitvijo na črnem in belem je definirano z $ratio$.

Opazim, da se rešitev razlikuje od homogene. Ker imamo tukaj nehomogeno porazdelitev mase, `zvezna metoda odpove'. Povesi so večji, kot pri homogeni in pričakovano malenkost večji pri središčni obtežbi. Zanimivo je, da z večanjem $ratio$ do neke mere spreminja oblika povesa, pri zelo velikih vrednostih pa oblika ostaja precej podobna, spremenijo se le vrednosti povesov. Ponovno povsod rešitev po MKR zaostaja, dočim po kombinirani metodi prehiteva direktno DFT.

\subsubsection{Obtežitev na robu}

\fig{H}{0.1}{./10.nal_slike/nehomogena-kvadrat na robu}{Opna z obtežitvijo na robu (črno)}{}
\begin{figure}[H]
    \centering
    \subfig{1}{./10.nal_slike/metode_nehomogena-kvadrat na robu_r2}{$ratio=2$}{}
    \subfig{1}{./10.nal_slike/metode_nehomogena-kvadrat na robu_r10}{$ratio=10$}{}
    \subfig{1}{./10.nal_slike/metode_nehomogena-kvadrat na robu_r100}{$ratio=100$}{}
    \caption{Poves nehomogene opne z različnimi razmerji obtežitve na robu}
\end{figure}
\begin{figure}[H]
    \centering
    \subfig{0.48}{./10.nal_slike/Umax_nehomogena-kvadrat na robu_r2}{$ratio=2$}{}
    \subfig{0.48}{./10.nal_slike/Umax_nehomogena-kvadrat na robu_r10}{$ratio=10$}{}
    \subfig{0.48}{./10.nal_slike/Umax_nehomogena-kvadrat na robu_r100}{$ratio=100$}{}
    \caption{Največji poves nehomogene opne z različnimi razmerji obtežitve na robu}
\end{figure}
\begin{figure}[H]
    \centering
    \subfig{1}{./10.nal_slike/prerez_nehomogena-kvadrat na robu_r2}{$ratio=2$}{}
    \subfig{01}{./10.nal_slike/prerez_nehomogena-kvadrat na robu_r10}{$ratio=10$}{}
    \subfig{1}{./10.nal_slike/prerez_nehomogena-kvadrat na robu_r100}{$ratio=100$}{}
    \caption{Prerez preko koordinatnih osi nehomogene opne z različnimi razmerji obtežitve na robu (črno)}
\end{figure}

\subsubsection{Obtežitev na sredini}

\fig{H}{0.1}{./10.nal_slike/nehomogena-kvadrat na sredini}{Opna z obtežitvijo na sredini (črno)}{}
\begin{figure}[H]
    \centering
    \subfig{1}{./10.nal_slike/metode_nehomogena-kvadrat na sredini_r2}{$ratio=2$}{}
    \subfig{1}{./10.nal_slike/metode_nehomogena-kvadrat na sredini_r10}{$ratio=10$}{}
    \subfig{1}{./10.nal_slike/metode_nehomogena-kvadrat na sredini_r100}{$ratio=100$}{}
    \caption{Poves nehomogene opne z različnimi razmerji obtežitve na robu}
\end{figure}
\begin{figure}[H]
    \centering
    \subfig{0.48}{./10.nal_slike/Umax_nehomogena-kvadrat na sredini_r2}{$ratio=2$}{}
    \subfig{0.48}{./10.nal_slike/Umax_nehomogena-kvadrat na sredini_r10}{$ratio=10$}{}
    \subfig{0.48}{./10.nal_slike/Umax_nehomogena-kvadrat na sredini_r100}{$ratio=100$}{}
    \caption{Največji poves nehomogene opne z različnimi razmerji obtežitve na robu}
\end{figure}
\begin{figure}[H]
    \centering
    \subfig{1}{./10.nal_slike/prerez_nehomogena-kvadrat na sredini_r2}{$ratio=2$}{}
    \subfig{1}{./10.nal_slike/prerez_nehomogena-kvadrat na sredini_r10}{$ratio=10$}{}
    \subfig{1}{./10.nal_slike/prerez_nehomogena-kvadrat na sredini_r100}{$ratio=100$}{}
    \caption{Prerez preko koordinatnih osi nehomogene opne z različnimi razmerji obtežitve na robu}
\end{figure}

\subsection{Obtežitev `f' iz naloge 6}

\fig{H}{0.1}{./10.nal_slike/f-opna}{`f˙opna iz 6.naloge}{}
Ponovno opazim presenetljivo značilnost, da $ratio$ pri majhnih vrednostih vpliva na obliko povesa, pri velikih pa na velikost povesa. To bi bilo vredno podrobneje raziskati.
\begin{figure}[H]
    \centering
    \subfig{1}{./10.nal_slike/metode_f-opna_r2}{$ratio=2$}{}
    \subfig{1}{./10.nal_slike/metode_f-opna_r10}{$ratio=10$}{}
    \subfig{1}{./10.nal_slike/metode_f-opna_r100}{$ratio=100$}{}
    \caption{Poves nehomogene opne z različnimi razmerji obtežitve}
\end{figure}
\begin{figure}[H]
    \centering
    \subfig{0.48}{./10.nal_slike/Umax_f-opna_r2}{$ratio=2$}{}
    \subfig{0.48}{./10.nal_slike/Umax_f-opna_r10}{$ratio=10$}{}
    \subfig{0.48}{./10.nal_slike/Umax_f-opna_r100}{$ratio=100$}{}
    \caption{Največji poves nehomogene opne z različnimi razmerji obtežitve}
\end{figure}
\begin{figure}[H]
    \centering
    \subfig{0.77}{./10.nal_slike/prerez_f-opna_r2}{$ratio=2$}{}
    \subfig{0.77}{./10.nal_slike/prerez_f-opna_r10}{$ratio=10$}{}
    \subfig{0.77}{./10.nal_slike/prerez_f-opna_r100}{$ratio=100$}{}
    \caption{Prerez preko koordinatnih osi nehomogene opne z različnimi razmerji obtežitve}
\end{figure}

\subsection{Konvergenca}

Raziskal sem še konvergenco metod in ugotovil, da prihaja do največjih razlik med SOR in posledično MRK in kombinirano metodo, ker vsaka konvergira iz svoje strani, dočim direktna DFT že pri malo računskih točkah vrne presenetljivo dober rezultat. Zadovolji me, da konvergirajo k isti vrednosti. Konvergenco sem meril glede vrednosti največjega povesa $u_{\max}$, vendar bi lahko to nadgradil z normo celotne rešitve.

Glede računske zahtevnosti sem meril čas in z linearno regresijo ugotovil približne rede konvergence. Pričakovano je pri SOR in zvezni občutno višji red, kot pri direktni DFT, ki je skoraj linearna. Zato pri sem se analizi z njima omejil na manjše število točk. Kombinirana metoda se presenetljivo obnese najbolje, 1D in 2D DFT sta enakovredni, vendar je 1D DFT malo počasnejša, verjetno zaradi večkratnega klica funkcij.

\begin{table}[H]
    \centering
    \begin{tabular}{c|c|c|c|c|c}
        metoda & 2D sin FFT & 1D sin FFT & kombinirana & zvezna & SOR \\ \hline
        $\mathcal{O}(N^\square)$ & 1,7 & 1,75 & 1,55 & 3,1 & 3,25
    \end{tabular}
    \caption{Ocenjeni red konvergence metod}
\end{table}
\begin{figure}[H]
    \centering
    \subfig{0.45}{./10.nal_slike/konvergenca/konvergenca_Umax_homog}{Homogena opna}{}
    \subfig{0.45}{./10.nal_slike/konvergenca/kvad_sred/konvergenca_Umax_N}{Kvadratna opna z obtežbo na sredini}{}
    \subfig{0.45}{./10.nal_slike/konvergenca/f_opna/konvergenca_Umax_N}{'f' opna}{}
    \caption{Konvergenca največjega povesa različnih obtežb}
\end{figure}
\begin{figure}[H]
    \centering
    \subfig{0.45}{./10.nal_slike/konvergenca/cas_homog}{Homogena opna}{} 
    \subfig{0.45}{./10.nal_slike/konvergenca/kvad_sred/cas_N}{Kvadratna opna z obtežbo na sredini}{}
    \subfig{0.45}{./10.nal_slike/konvergenca/f_opna/cas_N}{'f' opna}{}
    \caption{Računska zahtevnost}
\end{figure}

\section{Prevod toplote v valju}

\subsection{Temperaturni profil v valju}

V tem poglavju preizkusimo direktno metodo še v cilindričnem koordinatnem sistemu, kjer upoštevamo osno simetrijo, kar poenostavi Poissonovo enačbo,
\begin{equation}
    \frac{\partial^2 u}{\partial r^2} + \frac{1}{r} \frac{\partial u}{\partial r} + \frac{\partial^2 u}{\partial z^2} = g(r,\theta).
\end{equation}
Pri tem so robni pogoji pri $r=0$ Neumanovi in pri $r=R$, $z=0$ ter $z=L$ Dirichletovi. Za rešitev uporabim kombinirano metodo direktnega reševanja, kjer uporabim diskretno sinusno transformacijo (DST) v $z$ smeri in rešim tridiagonalni sistem enačb v $r$ smeri. Direktna DFT metoda v tem primeru odpove, saj sinusi niso bazne funkcije v radialni smeri. To so Bessli.

Na spodnji sliki so predstavljeni primer iz navodil in variacije s spremenjenimi robnimi pogoji na plašči valja. Pri tem se Neumanov robni pogoj na levi strani (notranjost valja) glasi
\begin{equation}
    Q_N=0,
\end{equation}
Dirichletovi robni pogoji pa so homogeni na zgornji in spodnji strani valja
\begin{equation}
    T_{D}=0.
\end{equation}
Spreminjamo samo Dirichletove robne pogoje na plašču valja:
\begin{itemize}
    \item konstantni nehomogeni RP na plašču (po navodilih) $T=1$,
    \item konstantni nehomogeni RP na zgornjem delu plašča $T=1$, sicer homogeno $T=0$,
    \item linearen RP na plašču $T(z)=z/L$,
    \item paraboličen RP na plašču $T(z)=\frac{4}{L}\;\left(z-\frac{z^2}{L}\right)$,
    \item sinusni RP na plašču $T(z)=\sin{\left(\frac{M\;\pi\;z}{L}\right)}$, kjer je $M=1$ za en sinusni val in $M=2$ za polni val.
\end{itemize}

Neumanove pogoje sem apliciral v tridiagonalnem sistemu enačb, kjer sem na glavno diagonalo dodal vrednost $1$ in na diagonalo nad njo vrednost $-1$, kar ustreza Neumanovemu pogoju. Na desni strani sem določil vrednost toplotnega toka preko roba $Q_N=0$. Homogenih Dirichletovih pogojev ni potrebno posebej obravnavati, saj so upoštevani že z izbiro DST kot transformacije. Nehomogene pogoje pa upoštevamo tako, da vrednosti na robovih dodamo v vektorju desnih strani kot dodatne izvire toplote, poleg volumske generacije toplote $Q_V$,
\begin{equation}
    g(r,z) = Q_V + \frac{T_D}{\Delta r^2}.
\end{equation}

Na spodnji sliki in v celotnem poglavju je vrednost nehomogenega Dirichletovega pogoja enaka $T_D=1$, na spodnji sliki so označeni še robni pogoji in izračunane povprečne vrednosti temperature na robu, kar je sicer smiselno le za konstantne RP kot merilo pravilnosti rešitve v smislu upoštevanja robnih pogojev.

\begin{figure}[H]
    \centering
    \subfig{0.3}{./10.nal_slike/valj/classic}{po navodilih -- konstantni nehomogeni RP na plašču}{}
    \subfig{0.3}{./10.nal_slike/valj/upper}{konstantni nehomogeni RP na zgornjem delu plašča}{}
    \subfig{0.3}{./10.nal_slike/valj/lin}{\\linearen}{}
    \subfig{0.3}{./10.nal_slike/valj/parabolic}{paraboličen}{}
    \subfig{0.3}{./10.nal_slike/valj/sine}{sinusni val}{}
    \subfig{0.3}{./10.nal_slike/valj/sine_full}{polni sinusni val}{}

    \caption{Različni Dirichletovi RP za potek temperature po plašču}
\end{figure}

Na spodnjih dveh slikah je prikazan še temperaturni profil za primer spremenjenega Neumanovega RP na sredini valja, kar sicer fizikalno ni smiselno, vendar je računsko precej enostavno, in volumskega izvora toplote v valju, kar je lažje fizikalno razumeti.

Pri tem opazimo, da lahko precej enostavno in učinkovito spreminjamo robne pogoje, tako Neumanove, kot tudi Dirichletove.

\begin{figure}[ht]
    \centering
    \subfig{0.3}{./10.nal_slike/valj/heat_source_2}{$Q_V=+2$}{}
    \subfig{0.3}{./10.nal_slike/valj/heat_source_5}{$Q_V=+5$}{}
    \subfig{0.3}{./10.nal_slike/valj/heat_source_-2}{$Q_V=-2$}{}
    \subfig{0.3}{./10.nal_slike/valj/heat_source_-5}{$Q_V=-5$}{}

    \caption{Volumski izvor toplote $Q_V$ v valju}
\end{figure}

\begin{figure}[H]
    \centering
    \subfig{0.3}{./10.nal_slike/valj/Q_BC2}{$Q_N=+2$}{}
    \subfig{0.3}{./10.nal_slike/valj/Q_BC-2}{$Q_N=-2$}{}

    \caption{Neumanov RP $Q_N$ na levi strani (notranjost valja)}
\end{figure}

\subsection{Konvergenca}

Ker nas zanima, kako točno metoda izpolnjuje robne pogoje, preverimo, kakšen je temperaturni profil na robu za različne diskretizacije. V tem poglavju obravnavamo problem iz navodil, kjer je na plašču konstanten Dirichletov nehomogen RP $T=1$. Pričakovano se z večanjem števila diskretnih točk izboljšuje točnost, vendar prihaja na robovih desno zgoraj in spodaj do diskretnega skoka temperature med 0 in 1, zaradi česar imamo temperaturo v vogalu okoli 0,5. Morebiti bi lahko vogalne točke izvzeli iz analize, saj očitno obstajajo težave pri apliciranju RP v njih, kar bi verjetno izboljšalo konvergenco. Prav tako sem zadovoljen, da se temperature na vogalih na levi strani skladajo in je potek na gornjem in spodnjem robu popolnoma enak.

\fig{H}{0.7}{./10.nal_slike/valj/conv_profil_RP}{Konvergenca temperaturnega profila in vizualno upoštevanje robnih pogojev}{}

Analitične rešitve nisem iskal, ampak sem rezultate primerjal z rešitvijo pri zelo fini diskretizaciji ($10^8$ računskih točk). Za primerjavo norm sem uporabil Frobeniusove norme, kar je prikazano na spodnjih dveh slikah.

Iz slike z relativno konvergenco norme opazim, da se konvergenca izboljšuje z večanjem števila diskretnih točk, vendar je precej počasna za, sploh za celoten temperaturni profil. Edino konvergenca na desni strani je očitno hitrejša. To si lahko razlagamo s tem, da je v primeru nehomogenega RP na desni strani pri manjšem številu diskretnih točk dosežena precej slaba točnost, ki pa se precej hitro izboljša z dodajanjem računskih točk. Temperatura se precej hitro približuje pričakovani vrednosti 1, kot je razvidno iz zgornje slike. Konvergence RP na sredini ($r=0$) in nehomogenih RP (zgoraj in spodaj) se prekrivajo.

\fig{ht}{0.6}{./10.nal_slike/valj/conv_norm_relat}{Konvergenca norme relativno glede na najgostejšo diskretizacijo}{}

Absolutna konvergenca je merjena z normo in deljena s številom računskih točk. Pri nehomogenih RP je odšteta še predpisana robna vrednost temperature. Trend je pričakovano padajoč in opazno hitrejši na desni strani, kjer je nehomogeni Dirichletov robni pogoj, kot na vsoti nehomogenih RP zgoraj in spodaj. To potrjuje, da se konvergenca izboljšuje z večanjem števila diskretnih točk.

\fig{ht}{0.6}{./10.nal_slike/valj/conv_norm_RP}{Absolutna konvergenca norme Dirichletovih RP glede na predpisano robno temperaturo}{}

\fig{ht}{0.6}{./10.nal_slike/valj/conv_time}{Časovna zahtevnost}{}

Iz zadnje slike s časovno zahtevnostjo je razvidno, da je krivulja skoraj linearna, če izvzamemo začetne nestabilnosti. Prav tako je naklon skoraj 45$^\circ$, od koder izhaja, da je red konvergence $\mathcal{O}(N) \approx 1$, kar je približno linearno in pričakovano od učinkovitega algoritma, kot je FFT, ki se skriva v ozadju DST (\texttt{scipy.fft.dst}).